% Copyright (c) 2026 Eugene Joseph M. Ragasa.
% SPDX-License-Identifier: Apache-2.0
% Author: Eugene Joseph M. Ragasa
% Formal proof completed: 2026-10-28
\documentclass[11pt]{article}
\usepackage[T1]{fontenc}\usepackage{amsmath,amssymb}\usepackage[margin=1in]{geometry}\usepackage[hidelinks]{hyperref}
\title{Proof Dossier: Matrix-Frobenius Fourier Truncation-Channel RMS}\author{Eugene Joseph M. Ragasa}\date{Formal proof completed: 2026-10-28}
\begin{document}\raggedright\maketitle
\section{Identity and status}
ID \texttt{BG-FOURIER-THM-072}; status \texttt{proved}. Lean declaration:
\par\noindent{\tiny\path{ProjectKoios.BandGauge.Fourier.finiteMatrixFrobeniusRootMeanSquare_truncationResidual_eq_sqrt_omitted}}.
\section{Statement and proof}
For complex matrix coefficients on a nonempty finite abelian group, the finite root-mean-square entrywise Frobenius norm of the deliberate-truncation channel equals the square root of the omitted coefficient Frobenius power. The proof applies the exact matrix Parseval tail identity and cancels the strictly positive square root of the group cardinality.
\section{Boundary and audit}
The RMS definition uses explicit entrywise Frobenius power, not the ambient matrix norm. No runtime matrix, mesh, or physical support interpretation is inferred. Axioms are \texttt{propext}, \texttt{Classical.choice}, and \texttt{Quot.sound}; no admissions. Lean \texttt{4.35.0}, Mathlib {\scriptsize\path{56a005855ea17e0b0958f73ec801ae140b0a1679}}; AI-assisted, kernel-checked, no independent expert review recorded.
\end{document}
